\section{拓扑环的完备化}

记号约定:$A$是拓扑环.

\begin{convention}{}{}
	论及拓扑环的一致结构时,我们总是指其加法群的一致结构.
\end{convention}

\begin{definition}{完备拓扑环}{}
	如果$A$的加法群是完备的,则称$A$是完备拓扑环.	
\end{definition}

\begin{theorem}{双线性映射的连续延拓}{}
	设$E,F,G$是Hausdorff的完备交换拓扑群,$P,Q$分别是$E,F$的稠密子群,$f\in\Bil_{\Z}\left(P,Q;G\right)$.若$f$连续,则$f$可以唯一地延拓为连续的双线性映射$\bar{f}:E\times F\to G$.
\end{theorem}

\begin{proposition}{Hausdorff拓扑环的完备化的存在性和唯一性}{existence-and-uniqueness-of-completion-of-hausdorff-topological-ring}
	设$A$是Hausdorff拓扑环,则存在完备Hausdorff拓扑环$\hat{A}$和$\iota\in\Hom_{\mathsf{TopRing}}(A,\hat{A})$满足如下条件:
	\begin{enumerate}
		\item 对每个完备Hausdorff拓扑环$B$和$f\in\Hom_{\mathsf{TopRing}}\left(A,B\right)$,存在唯一的$\bar{f}\in\Hom_{\mathsf{TopRing}}(\hat{A},B)$使下图交换.
		\begin{center}
			\begin{tikzcd}
				A && B \\
				& {\hat{A}}
				\arrow["f", from=1-1, to=1-3]
				\arrow["\iota"', from=1-1, to=2-2]
				\arrow["{\exists!\bar{f}}"', dashed, from=2-2, to=1-3]
			\end{tikzcd}
		\end{center}
		\item 对每个满足(1)的每份资料$\left(A^{\prime},\iota^{\prime}\right)$,存在唯一的$\bar{\pi}\in\Isom_{\mathsf{TopRing}}(\hat{A},A^{\prime})$使下图交换.
		\begin{center}
			\begin{tikzcd}
				A && {A^{\prime}} \\
				& {\hat{A}}
				\arrow["{\iota^{\prime}}", from=1-1, to=1-3]
				\arrow["\iota"', from=1-1, to=2-2]
				\arrow["{\exists!\bar{\pi}}"', dashed, from=2-2, to=1-3]
			\end{tikzcd}
		\end{center}
	\end{enumerate}
\end{proposition}

\begin{remark}{}{}
	如果$A$是交换环,那么$A^{\prime}$也是交换环;如果$A$有单位元,那么$A^{\prime}$也有单位元.
\end{remark}

\begin{definition}{Hausdorff拓扑环的完备化}{}
	设$A$是Hausdorff拓扑环.我们称\cref{prop:existence-and-uniqueness-of-completion-of-hausdorff-topological-ring}中的$\hat{A}$是$A$的Hausdorff完备化,称$\iota$是完备化映射.
\end{definition}

\begin{definition}{拓扑环的Hausdorff完备化}{}
	设$N=\closure\left(\left\{0\right\}\right),A^{\prime}=A/N$.我们把$A^{\prime}$的完备化记作$\hat{A}$,称之为$A$的Hausdorff完备化.
\end{definition}

\begin{proposition}{Hausdorff完备化的泛性质}{}
	设$B$是拓扑环,$u\in\Hom_{\mathsf{TopRing}}\left(A,B\right)$,则存在唯一的$\hat{u}\in\Hom_{\mathsf{TopRing}}(\hat{A},\hat{B})$使下图交换(其中$\iota,\iota^{\prime}$是完备化映射).
	\begin{center}
		\begin{tikzcd}
			A & B \\
			{\hat{A}} & {\hat{B}}
			\arrow["u", from=1-1, to=1-2]
			\arrow["\iota"', from=1-1, to=2-1]
			\arrow["{\iota^{\prime}}", from=1-2, to=2-2]
			\arrow["{\exists!\hat{u}}"', dashed, from=2-1, to=2-2]
		\end{tikzcd}
	\end{center}
\end{proposition}
